\chapter{Getting Started Together}
\label{ch:orientation}

\section*{Let me tell you how I approach a data project}

If we were sitting down together with these files for the first time, I would
not begin by opening SQL or trying to build a model. I would begin with a few
ordinary questions: Where did these records come from? Who created them? What
does one row describe? What might be missing? What would I need to check before
I trusted a number?

That is how I begin my work as a data scientist. I first learn the story of the
data. Later, you will see technical names for parts of that process, but the
idea itself is simple: before we calculate anything, we should understand what
we have in front of us.

In this project, we have seven files about courses, students, registrations,
online activity, and assessments. We are going to open them carefully, move
them into a database, ask questions with SQL, and check our answers as we go.
I will explain the choices I made in the project, including the places where I
had to slow down or be cautious.

You do not need to memorize every new word on the first pass. I will introduce
the language when it becomes useful, and we will keep returning to the same
central question: \emph{What does the data actually allow us to say?}

\begin{learningobjectives}
After completing this chapter, you will be able to:
\begin{itemize}
  \item state what you will build in Course Edition 0.1;
  \item distinguish project implementation from course documentation;
  \item identify the prerequisites, checkpoints, and evidence standards;
  \item describe the complete project workflow at a high level; and
  \item use the course's read-run-validate-reflect study cycle.
\end{itemize}
\end{learningobjectives}

\section{A few words we will use together}

\begin{plainlanguage}
This is a translation guide, not a test. At the simplest level, the project is
a folder with data files that hold recorded information and instruction files
that tell the computer what to do with it. The terms below give us shorter
names for those familiar ideas.
\end{plainlanguage}

The terms in \cref{tab:first-vocabulary} will appear throughout the course.
Skim them now and return whenever a word feels unfamiliar. I do not expect you
to memorize the table before continuing.

\begin{longtable}{@{}L{0.22\textwidth}L{0.71\textwidth}@{}}
\caption{First vocabulary for the course}\label{tab:first-vocabulary}\\
\toprule
Term & Meaning in this course \\
\midrule
\endfirsthead
\toprule
Term & Meaning in this course \\
\midrule
\endhead
\term{Data} &
Recorded facts or values. A \term{dataset} is an organized collection of those
records. Data can be useful without being complete or perfectly accurate. \\
\term{Data science} &
The work of turning recorded information into an answer I can explain and
stand behind. Programming, databases, and statistics help me do that work, but
the work begins with understanding the data and the question. \\
\term{File, folder, path} &
A file stores content. A folder, also called a directory, groups files. A path
is the address of a file or folder, such as
\projectfile{sql/02_staging.sql}. \\
\term{CSV file} &
A plain-text table in which values are separated by commas. CSV means
comma-separated values. Chapter 3 opens one and explains every part. \\
\term{Archive} &
One packaged file that contains other files, commonly in ZIP format. Projects
often download one archive and then unpack its contents. \\
\term{Table, row, column} &
A table arranges data in a grid. A row usually describes one recorded thing or
event. A column contains one kind of value, such as a student identifier or
date. \\
\term{Code and script} &
Code is text written as instructions for a computer. A script is a file of
code intended to perform a sequence of steps. \\
\term{Terminal and command} &
A terminal is a text-based way to interact with the computer. A command is one
instruction typed there, such as asking PostgreSQL to run a SQL file. \\
\term{Database} &
Software-managed storage for organized data. Unlike a single spreadsheet, a
database can hold many related tables and enforce rules about them. \\
\term{PostgreSQL} &
The database program used by this project. It stores the OULAD tables and runs
the project's SQL instructions. \\
\term{SQL and query} &
SQL is a language for working with relational databases. A query is a SQL
request, often asking PostgreSQL to return, combine, summarize, or check rows. \\
\term{Python} &
A general-purpose programming language. This project uses Python to download
and audit files and, in later stages, to analyze, chart, and model data. \\
\term{Repository} &
The main project folder and its organized contents. The word often implies
that file history is tracked with Git, a file-history tool, but you do not need
Git to read the course or query the database. \\
\term{Implementation and documentation} &
The implementation is the working code and settings that perform the project.
Documentation explains what that implementation does and how to use or
evaluate it. \\
\term{Input and output} &
An input is what a step receives. An output is what it produces. The official
CSV files are inputs; checked tables and charts are outputs. \\
\term{Workflow or pipeline} &
An ordered sequence in which the output of one step becomes the input to the
next. \emph{Pipeline} is the common technical word; \emph{workflow} is the same
idea in plainer language. \\
\term{Analysis, model, dashboard} &
Analysis examines data to answer a question. A model is a mathematical rule
learned or estimated from data. A dashboard is a visual display of selected
results. They are different products and should not be confused. \\
\term{Reproducible} &
A result is reproducible when another person can follow the recorded source,
code, settings, and steps and obtain the same result. \\
\bottomrule
\end{longtable}

\begin{validationbox}
Before continuing, make sure you can explain this sentence:
\begin{quote}
``A script reads an input CSV file, loads its rows into PostgreSQL, and a SQL
query produces an output table.''
\end{quote}
It is fine to return to the vocabulary table while translating it into your own
words.
\end{validationbox}

\section{The project we will work through together}

I built this course around a real data project:
\emph{Learning Analytics: Engagement, Assessment, and Outcome Forecasting}.
The project uses the Open University Learning Analytics Dataset (OULAD), a set
of seven linked CSV files published for research and education
\citep{oulad_dataset,oulad_paper}.

Here is the path we will follow:

\begin{figure}[htbp]
\centering
\begin{tikzpicture}[
  node distance=6mm,
  box/.style={
    draw=CourseBlue,
    rounded corners=1mm,
    fill=CourseLight,
    minimum width=0.82\textwidth,
    minimum height=8mm,
    align=center,
    font=\small
  },
  arrow/.style={-{Latex[length=2.4mm]},thick,color=CourseGold}
]
\node[box] (source) {Official OULAD CSV data files};
\node[box,below=of source] (audit) {Check that the downloaded files are complete and unchanged};
\node[box,below=of audit] (database) {Store the original values, then convert them into usable tables};
\node[box,below=of database] (quality) {Test the tables and prepare carefully defined summaries};
\node[box,below=of quality] (python) {Analyze, chart, and eventually make and test predictions};
\node[box,below=of python] (output) {Save tables, figures, findings, and a results dashboard};
\draw[arrow] (source) -- (audit);
\draw[arrow] (audit) -- (database);
\draw[arrow] (database) -- (quality);
\draw[arrow] (quality) -- (python);
\draw[arrow] (python) -- (output);
\end{tikzpicture}
\caption{The path we will follow through the project. Edition 0.1 takes us
through the SQL foundations stage; later lessons will continue the work.}
\label{fig:project-lineage}
\end{figure}

\subsection{The shorter names used in the project files}

The diagram begins in everyday language. When you open the project files, you
will see the following shorter names for the same steps:
\begin{description}
  \item[Checksum] A compact fingerprint used to check whether a file's bytes
    match the expected file.
  \item[Audit] A recorded inspection of source files, such as checking names,
    row counts, columns, and missing values.
  \item[Raw layer] Database tables that preserve source text as closely as
    practical.
  \item[Staging layer] An intermediate area where text is converted into
    meaningful types, names, and missing values.
  \item[Core layer] The main checked tables used for analysis.
  \item[Data validation] Tests that compare the data with explicit
    expectations.
  \item[Snapshot] A saved view of information as of a particular course week.
  \item[Lineage] The traceable path from original input to a reported result.
\end{description}

\begin{plainlanguage}
We are not starting with an empty practice exercise. We are taking apart a
working project together. The code shows what the computer does; my job in this
course is to explain why the steps are there, what choices they contain, and
how we can check them.
\end{plainlanguage}

\section{Edition 0.1 boundary}

The first edition fully develops the foundation shown in
\cref{tab:edition-scope}. It deliberately does not rush through the modeling
half of the project. Those chapters need their own worked examples, exercises,
and careful treatment of temporal leakage, which means accidentally using
information from the future to make an earlier prediction.

\begin{table}[htbp]
\centering
\caption{Course Edition 0.1 scope}
\label{tab:edition-scope}
\begin{tabularx}{\textwidth}{L{0.27\textwidth}YY}
\toprule
Area & Status in this edition & Concrete learner output \\
\midrule
Learning analytics and OULAD & Complete & Careful question statement and source-table map \\
Relational foundations & Complete & Row meaning (grain), key, and relationship worksheet \\
PostgreSQL setup & Complete & Verified database connection \\
Raw import & Complete & Seven populated raw tables \\
Data validation & Complete & Interpreted quality-check report \\
Introductory SQL & Complete & Reproducible OULAD query set \\
Joins through modeling & Forthcoming & Architecture only; no claim of completion \\
\bottomrule
\end{tabularx}
\end{table}

\section{Who I wrote this for}

I wrote this for someone who wants to learn SQL and data work by staying with
one project long enough to understand it. You do \emph{not} need prior OULAD
knowledge or advanced SQL. We will build the database and SQL ideas as we need
them.

\subsection{Prerequisites}

\begin{longtable}{L{0.27\textwidth}L{0.28\textwidth}L{0.36\textwidth}}
\caption{Prerequisite expectations and support}\label{tab:prerequisites}\\
\toprule
Skill & Expected starting point & Support in this edition \\
\midrule
\endfirsthead
\toprule
Skill & Expected starting point & Support in this edition \\
\midrule
\endhead
Files and folders & Able to locate a project folder & Exact paths from the main project folder are supplied \\
Terminal & Able to run a copied command & Commands are explained before use \\
Python & No pandas experience required yet & Installation is covered; Python analysis is forthcoming \\
PostgreSQL & No prior database administration required & Two supported setup paths are taught \\
SQL & No prior SQL required & Starts with \texttt{SELECT}, not advanced windows \\
Statistics & Basic percentages are helpful & Formal statistical reasoning is forthcoming \\
Git & Useful but not required for the lessons & No Git operation is required to query the database \\
\bottomrule
\end{longtable}

\section{How I suggest we work through each chapter}

I suggest that we use the same five-step rhythm throughout the course:

\begin{enumerate}
  \item \textbf{Read the question.} Begin with the analytical reason for a
    command or query.
  \item \textbf{Predict the output.} State the expected row grain, columns, and
    possible missing values before running code. The \term{grain} is what one
    output row will represent.
  \item \textbf{Run the smallest useful step.} Do not execute the entire
    workflow merely to learn one SQL clause, or instruction.
  \item \textbf{Validate.} Compare counts, keys, ranges, and categories with a
    known contract.
  \item \textbf{Reflect and transfer.} Explain what would change with a
    different dataset.
\end{enumerate}

\begin{projectconnection}
The project itself follows this logic. For example,
\projectfile{sql/04_quality.sql} does not merely calculate summaries. It records
the expected rules, called a data contract, counts rows that break them, records
how serious each failure is, and preserves the result for later inspection.
\end{projectconnection}

\section{A map of the repository}

\begin{table}[htbp]
\centering
\caption{Project paths used in Edition 0.1}
\label{tab:repo-map}
\begin{tabularx}{\textwidth}{L{0.35\textwidth}Y}
\toprule
Path & Role in the learning sequence \\
\midrule
\projectfile{README.md} & Project scope, architecture, results, and reproduction order \\
\projectfile{DATA_LICENSE.md} & Dataset citation, license, and access note \\
\projectfile{src/learning_analytics/download.py} & Download and file-fingerprint verification \\
\projectfile{src/learning_analytics/audit.py} & Source profiling and saved audit outputs \\
\projectfile{sql/00_database.sql} & Six named database areas, called schemas \\
\projectfile{sql/01_raw_tables.sql} & Source-faithful landing tables \\
\projectfile{sql/02_staging.sql} & Types, names, and missing-value normalization \\
\projectfile{sql/03_core.sql} & Keys, constraints, and core relational tables \\
\projectfile{sql/04_quality.sql} & Executable data-quality framework \\
\projectfile{sql/gallery/queries.sql} & Query sequence used as the SQL source bank \\
\projectfile{reports/tables/} & Verified source counts and quality results \\
\bottomrule
\end{tabularx}
\end{table}

\section{How we will keep the work honest}

I do not want you to trust a number just because it looks polished. In this
course, every numeric claim must lead back to a saved result or a query we can
run. For example, the statement ``OULAD contains 10,655,280 VLE activity
records'' is supported by
\projectfile{reports/tables/source_file_profile.csv} and a SQL row-count check.
If I show a teaching query whose result has not been saved, I will describe
what its columns and rows should look like instead of inventing a result.

\section{Knowledge check}

\begin{enumerate}
  \item Which stages of \cref{fig:project-lineage} are fully taught in Edition
    0.1?
  \item What is the difference between a verified result and an instructional
    query?
  \item Why does the course ask you to predict output grain before running SQL?
  \item Which skills are legitimate prerequisites, and which are taught here?
\end{enumerate}

\begin{exercisebox}{Exercise 1.1: Write your study contract}
In four short paragraphs, record:
\begin{enumerate}
  \item what you expect to build by the end of Edition 0.1;
  \item which prerequisite is currently weakest;
  \item how you will validate a result before trusting it; and
  \item which forthcoming topic you most want to reach.
\end{enumerate}
Your output should be a one-page study note. No code is required.
\solutionref{sol:orientation}
\end{exercisebox}

\begin{commonmistakes}
\begin{itemize}
  \item \textbf{Treating the dashboard as the course.} It communicates results;
    it does not teach the implementation.
  \item \textbf{Running everything immediately.} Large imports obscure basic
    concepts. Follow the chapter checkpoints.
  \item \textbf{Equating completion with understanding.} A command that exits
    successfully can still produce rows at the wrong level of detail.
\end{itemize}
\end{commonmistakes}

\transferheading
Before working with any unfamiliar project folder, draw the path from source to
output, distinguish generated files from editable source, and decide which
claims can be regenerated. This habit transfers to nearly every analytical
project.

\begin{summarybox}
Edition 0.1 is a bounded, project-based course release. It teaches the OULAD
context, the organization of linked database tables, PostgreSQL setup, source
import, validation, and SQL foundations. It does not pretend that the later
analytical and modeling chapters are finished. Each lesson follows a read,
predict, run, validate, and reflect cycle.
\end{summarybox}
